% chktex-file 13
\section{Introduction}
As trains become automated, the need for operators to monitor and control them remotely is growing. To do this safely, the operator needs a live video feed from the train that is stable and fast enough to react to what is happening on the track. This thesis looks at how such a video streaming system can be built, tested and evaluated, with a focus on finding configurations that work well for remote train operation.

\subsection{Background and Motivation}
As the railway sector transitions toward higher Grades of Automation (GoA), particularly GoA 3 and GoA 4, remote driving becomes an increasingly critical operational capability. Remote driving is not only a natural consequence of reduced onboard personnel at higher GoA levels, but serves as a mandatory safety layer within the ATO architecture enabling operators to intervene when automated systems encounter situations they cannot resolve independently, such as unexpected track conditions or degraded operational modes~\cite{orr2024, europesrail2025_d3}.

A fundamental requirement for safe remote driving is the availability of reliable, low-latency video streaming. When the operator is physically separated from the train, the video feed becomes the primary source of situational awareness, and its quality and responsiveness directly determine whether timely and safe intervention is possible~\cite{mejias2024, myhrvold2025}. Industry experience further demonstrates that remote operation can support a range of practical tasks including shunting, depot movements, and vehicle preparation provided that the communication link meets sufficient performance criteria~\cite{fp2r2dato2023d5_4}.

Understanding the relationship between latency and human perception is equally important in this context. Research on human factors in remote operation indicates that video delay, bitrate, and stream stability all influence an operator's ability to detect obstacles, interpret dynamic scenes, and execute corrective actions under time pressure~\cite{brandenburger2023, neumeier2019}. Identifying the latency boundaries where safe remote operation remains usable is therefore a key step toward enabling higher automation levels in railway systems~\cite{myhrvold2025}.

This thesis is developed within an ongoing research collaboration at NTNU, in partnership with VTI and other European railway organizations, contributing to the technical groundwork required to make remote train driving a safe and practical reality in future railway deployments.

\subsection{Problem}

The current remote train operation setup relies on a commercial product, Voysys, which imposes significant limitations on the project's ability to test and develop the system further. As a closed external solution, it restricts access to the underlying pipeline, making it difficult to adjust configurations, measure internal latency, or experiment with different streaming approaches. In a research context where iterative testing and configuration control are essential, this represents a fundamental obstacle. In addition, licensing such a product over the course of a long term research project represents a considerable and ongoing cost.

At the same time, there is limited existing research that addresses video streaming specifically in the context of remote train operation. Most work on low latency video transmission is either too general or focused on other domains, leaving a gap when it comes to understanding what configurations and compromises are most suitable for railway use cases [source].

Building a custom streaming solution also comes with its own set of challenges. Network conditions such as signal strength, interference, and varying load introduce unpredictable variability in transmission quality. Choosing the right codec, transport protocol, and encoder settings requires systematic testing, as each parameter affects latency, quality, and resource usage in ways that are not always predictable. Balancing video quality against latency is a difficult comparison and testing these differenes in settings on real hardware in a realistic setup adds further complexity. Without a flexible and open pipeline, it is difficult to isolate and measure where latency is introduced and how different configurations affect overall system performance.

\subsection{Objectives}

The main objective of this thesis is to design and implement a low-latency video streaming pipeline tailored specifically for remote train operation, serving as a practical replacement for the current commercial solution. The pipeline should be functional, flexible, and open enough to support continued development and testing within the project.
More specifically, the objectives are to:

\begin{enumerate}
    \item O1: Develop a working GStreamer-based streaming pipeline that meets the low-latency requirements of remote train operation
    \item O2: Implement a latency measurement framework that allows end-to-end delay to be captured, logged, and compared across different configurations
    \item O3: Evaluate the pipeline across multiple configurations, including different codecs, transport protocols and network settings, to identify the interplay between latency and video quality
    \item O4: Design the system to be dynamic and extensible, allowing additional cameras and new configurations to be added with minimal effort
    \item O5: Deliver a stable and usable prototype that is good enough for practical use during the current development period and serves as a foundation for future work
\end{enumerate}

The result should be a system that gives the project full control over the streaming pipeline, removes dependency on external licensed software, and produces reliable data to guide future configuration decisions.

\subsection{Research Questions}
The following research questions guide the experimental work of this thesis:
\begin{itemize}
    \item RQ1: What end-to-end latency and video quality can be achieved in a GStreamer-based video streaming pipeline for remote train operation, and which codec and transport configuration performs best?
    \item RQ2: How do network conditions, including transport protocol and available bandwidth, affect streaming latency and reliability in a remote train operation context?
    \item RQ3: What are the practical interplay between latency, video quality, and resource usage across different streaming configurations, and which configuration is most suitable for remote train operation?
\end{itemize}

\subsection{Scope and Limitations}

The scope of this thesis covers the design, implementation and experimental evaluation of a custom video streaming pipeline for remote train operation. The system is developed as a proof of concept and is not intended as a deployment ready solution. Development and pre testing was conducted in a controlled indoor environment, while final testing was performed from a moving vehicle to simulate realistic train movement conditions.

The experimental evaluation focuses on the parameters identified as most relevant to remote train operation based on prior research [source: prosjektoppgave], and selected based on time and hardware availability. The configurations tested include codec selection (H.264, H.265 and MJPEG), transport protocol (UDP and TCP), encryption overhead (SRTP), bandwidth conditions, resolution, and GPU versus CPU based encoding. While GStreamer allows a large number of additional parameters to be adjusted, such as buffer sizes, queue configurations, bitrate control modes and keyframe intervals, expanding the test matrix further was not feasible within the time constraints of a single master thesis. The chosen parameters are considered sufficient to identify the most significant performance differences relevant to the use case.

Latency is measured end to end between the point where an RTP packet enters the sender pipeline and the point where the corresponding decoded frame arrives at the receiver display sink. Camera capture latency and screen rendering latency are not included in these measurements, as they occur outside the instrumented portion of the pipeline and are considered largely constant across configurations. Clock synchronization between the sender and receiver relies on internet based time sources, introducing a small margin of uncertainty estimated at a few milliseconds, which is acknowledged when interpreting absolute latency values.

The thesis does not include testing on an actual operational train. A moving vehicle was used to simulate train conditions, but real world railway infrastructure, operational speeds, and trackside network coverage are outside the scope of this work and represent a natural next step for future validation.

The thesis does not include formal safety certification or a RAMS analysis in accordance with EN 50126. While the system is developed with safety critical use in mind, evaluating compliance with railway safety standards is considered outside the scope of this work and is left for future development phases.

\begin{comment}
With RAMS:
The system is developed with reference to EN 50126, the European railway standard for reliability, availability, maintainability and safety. While a full RAMS analysis is beyond the scope of a single master thesis, relevant safety considerations from the standard are used to inform design decisions and contextualize the results within a safety critical operational environment.
\end{comment}

\subsection{Research Methodology}
\subsubsection{Literature and Knowledge acquisition}
The knowledge needed to carry out this thesis was acquired through an iterative process that combined systematic literature review with hands-on practical exploration. This approach, often referred to as an Iterative Literature and Practice-Driven Approach, is characterized by a continuous cycle where theory and practice inform each other rather than following a strict linear sequence [Wohlin et al. (2012)]. Instead of completing all literature review before beginning practical work, the two activities run in parallel, with new practical challenges prompting targeted literature searches and new theoretical insights shaping the direction of experiments. This makes it particularly well suited for applied engineering research where the problem space is not fully known at the outset and understanding develops gradually through doing [Hevner et al. (2004) Saunders et al. (2019) ].

The process began by identifying the core knowledge areas required before any development could start. This included understanding the operational context of remote train operation and Automatic Train Operation, the fundamentals of video streaming and relevant transport protocols, and the capabilities and structure of GStreamer as a framework. These areas formed the initial scope of the literature review and gave enough grounding to begin practical work.

Academic literature was gathered primarily through Google Scholar and IEEE Xplore, using search terms related to low-latency video streaming, remote train operation, GStreamer pipeline design, and network performance for real-time applications. Sources were selected based on relevance to the problem, publication quality, and recency where applicable. In addition to academic papers, official technical documentation was used extensively, particularly the GStreamer documentation and protocol specifications for RTP and UDP, which provided the practical detail that academic sources alone could not cover.

Once a sufficient theoretical foundation was in place, practical exploration began through hands-on work with GStreamer. Building and testing pipelines from the command line and through Python quickly revealed challenges that existing literature did not fully address, such as specific configuration trade-offs, unexpected latency behavior, and hardware compatibility issues. Each of these practical obstacles prompted a return to the literature, this time with more precise and targeted questions. New findings were then brought back into the development process, tested, and evaluated against real behaviour.

This cycle of reading, building, observing, and searching again repeated throughout the entire development phase. The result is that the theoretical foundation of this thesis did not precede the practical work, but grew alongside it, shaped by the specific demands of the problem at hand.



\subsubsection{Search Strategy and Databases}
A multi-platform search approach was utilized to retrieve a wide array of high-quality sources. The primary databases included:

\begin{itemize}
    \item Scopus: For retrieving peer-reviewed, impactful scientific articles.
    \item Google Scholar: For broader academic and institutional literature.
    \item GStreamer docs: For specifics on GStreamer development.
\end{itemize}



\subsubsection{Inclusion Criteria}
To ensure the study is based on the relevant and current information, specific criteria were applied to filter the search results. Given the rapid pace of technological change, a focus was placed on recent publications. Where modern papers have referenced to older papers, I have include some as to say and show where certain numbers are coming from and how they have influenced results.

\begin{table}[h]
    \centering
    \caption{Inclusion Criteria for Literature Review}\label{tab:inclusion_criteria}
    \begin{tabular}{|c|p{8cm}|}
        \hline
        ID & Inclusion Criteria (IC) \\
        \hline
        IC1 & Publication Date: Between 2020-2026 \\
        \hline
        IC2 & Language: Written in English \\
        \hline
        IC3 & Document Type: Primarily "Article", "Conference Paper" or "Thesis"\\
        \hline
    \end{tabular}
\end{table}


\subsubsection{Search Queries and Refinement}

Several digital tools as mentioned below were used for this paper. The main purpose of the usage was to enhance clarity, assisting in latex, and in general increase quality. Specifics include writing reference list in correct format. Structuring sentences and paragraphs. Checking for typos and grammatical errors.

\begin{itemize}
    \item OpenAI's ChatGPT: Employed as a helpful resource for LaTeX formatting suggestions, generating structured content (tables and lists), translating between Norwegian and English, and reviewing text for synonyms and restructuring ideas to ensure arguments were effectively communicated.
    \item Google Gemini: Used for similar fields as ChatGPT, but also for general proofreading, and refining sentence structure and tone to maintain a high standard of academic writing.
    \item Anthropic's Claude.ai for discussion and development of script syntax.
    \item Visual Studio Code: was used as the main LaTeX editor because of its powerful extensions for LaTeX support, syntax highlighting, and version control integration.
\end{itemize}

\subsection{Outline of Structure}
The remainder of this thesis is organized as follows:

Chapter 2 presents the theoretical background underlying the work. This includes an overview of Automatic Train Operation and the role of remote driving at higher GoA levels, fundamentals of video streaming including codecs, transport protocols and payloading, an introduction to GStreamer and its core concepts, the main sources of latency in a streaming pipeline, relevant aspects of cellular and WiFi communication, and the metrics used to evaluate streaming performance.

Chapter 3 describes the design of the streaming system. This covers the requirements that shaped the system, the overall architecture and component choices, the GStreamer pipeline design for both sender and receiver, the approach to timestamp based latency measurement, and the network setup used in testing.

Chapter 4 explains how the research was conducted. This includes the overall research design, the iterative development process and how it shaped the system over time, the GStreamer pipeline development process, and the experimental and latency measurement methodology used to evaluate the system.

Chapter 5 presents the experimental results. This covers a comparison of GStreamer configurations across different codecs, transport protocols and encoding approaches, latency measurements per configuration, the effect of network conditions on streaming performance, and the interplay between video quality and latency across configurations.

Chapter 6 interprets and reflects on the results in relation to the research questions and the broader context of remote train operation.

Chapter 7 summarizes the key findings and contributions of the thesis and outlines directions for future work, including testing on real railway infrastructure and further expansion of the configuration space.